\subsection{例}

在本小节中（例 7 和 8 除外），K 表示一个非离散的局部紧交换域；dx 表示 K 的加法群上的一个 Haar 测度。

回忆：\(\mod x = |x|\) 当 \(\mathbf{K} = \mathbf{R}\) 时，\(\mod x = |x|^2\) 当 \(\mathbf{K} = \mathbf{C}\) 时，\(\mod x = |x|_p\) 当 \(\mathbf{K} = \mathbf{Q}_p\) 时。

例 1. --- 一般线性群。

令 A 为代数 \(\mathbf{M}_{n}(\mathbf{K})\)。A 的可逆元群 \(A^{*}\) 正是一般线性群 \(\mathbf{GL}(n,\mathbf{K})\)。对于每个 \(X\in A\)，约化范数 \(\mathrm{Nrd}_{\mathrm{A}/\mathrm{K}}(X)\) 是 det X；因此 \(\mathrm{N}_{\mathrm{A}/\mathrm{K}}(X)=(\det X)^{n}\)（Alg., Ch. VIII, §12, No.~3, Prop. 8；cf.~A, III, §9, No.~3, Example 3）。由于 \(X\mapsto{}^{t}X\) 是 A 到其逆代数上的同构，

\[
\mathrm{N} _ {\mathrm{A} ^ {0} / \mathrm{K}} (X) = \mathrm{N} _ {\mathrm{A} / \mathrm{K}} (^ {t} X) = \det (^ {t} X) ^ {n} = (\det X) ^ {n}.
\]

于是，§1, No.~11 的 Prop. 16 表明，测度

\[
\operatorname{mod} (\det X) ^ {- n} \cdot \bigotimes_ {i, j} d x _ {i j} \quad (X = (x _ {i j}))\tag{3}
\]

是 \(\mathbf{GL}(n,\mathbf{K})\) 上的左、右 Haar 测度。

为确定 \(\mathbf{GL}(n,\mathbf{K})\) 上的相对不变测度，我们将依靠以下引理：

引理 5. --- \(\mathbf{GL}(n,\mathbf{K})\) 到 \(C^{*}\) 的连续表示正是形如 \(X \mapsto \chi(\det X)\) 的映射，其中 \(\chi\) 是 \(K^{*}\) 到 \(C^{*}\) 的连续表示。

这样的映射显然是 \(\mathbf{GL}(n,\mathbf{K})\) 到 \(C^{*}\) 的连续表示。反过来，设 \(\psi\) 是 \(\mathbf{GL}(n,\mathbf{K})\) 到 \(C^{*}\) 的连续表示。对 \(x \in K^{*}\)，令

\[
\widetilde {x} = \left( \begin{array}{c c c c c} x & & & & \\ & 1 & & 0 & \\ & & 1 & & \\ & 0 & & \ddots & \\ & & & & 1 \end{array} \right)
\]

并令 \(\chi(x) = \psi(\widetilde{x})\)。于是，对于每个矩阵 \(X \in \mathbf{GL}(n, K)\)，有 \((\det X^{-1})^{\sim} \cdot X \in \mathbf{SL}(n, K)\)。由于 \(\mathbf{SL}(n, K)\) 是 \(\mathbf{GL}(n, K)\) 的换位子群（A, III, §8, No.~9, Prop. 17 的 Cor.），\(\psi((\det X^{-1})^{\sim} \cdot X) = 1\)，从而

\[
\psi (X) = \psi \big ((\det X) ^ {\sim} \big) = \chi (\det X).
\]

由此，§1, No.~8 的 Prop. 10 的 Cor. 1 表明，\(\mathbf{GL}(n,\mathbf{K})\) 上的相对不变测度除一个常数因子外，都是如下形式的测度

\[
\chi (\det X) \cdot \bigotimes_ {i j} d x _ {i j} \quad (X = (x _ {i j})),\tag{4}
\]

其中 \(\chi\) 是 \(\mathbf{K}^*\) 到 \(\mathbf{C}^*\) 的连续表示。

例 2. --- 仿射群。

对于每个 \(X \in \mathbf{GL}(n, \mathbf{K})\) 和每个 \(x \in K^{n}\)，令 \((X, x)\) 为 \(\xi \mapsto X\xi + x\) 在 \(K^{n}\) 上的仿射线性映射。所有 \((X, x)\) 构成 \(K^{n}\) 的仿射群 G（A, II, §9, No.~4）。平移集 T 是 G 的闭正规子群，典范地同构于 \(K^{n}\)；另一方面，\(\mathbf{GL}(n, \mathbf{K})\) 是 G 的闭子群，而 G 是 \(\mathbf{GL}(n, \mathbf{K})\) 与 \(T = K^{n}\) 的半直积。赋予 G 以这样的（局部紧）拓扑，使得 G 成为 \(\mathbf{GL}(n, \mathbf{K})\) 与 T 的拓扑半直积（GT, III, §2, No.~10）。有

\[
(X, x) = (1, x) \cdot (X, 0).
\]

另一方面，若 \(X \in \mathbf{GL}(n, \mathbf{K})\) 且 \(x \in \mathbf{T}\)，则对于每个 \(\xi \in \mathbf{K}^n\)，

\[
(X, 0) (1, x) (X, 0) ^ {- 1} \xi = X (X ^ {- 1} \xi + x) = \xi + X x = (1, X x) \xi ,
\]

因此，T 的自同构 \((1,x)\mapsto(X,0)(1,x)(X,0)^{-1}\) 的模为 mod(det X)（§1, No.~10, Prop. 15）。根据例 1 以及 §2, No.~9 的备注，测度

\[
\operatorname{mod} (\det X) ^ {- n - 1} \cdot \left(\bigotimes_ {i j} d x _ {i j}\right) \otimes \left(\bigotimes_ {i} d x _ {i}\right) \quad (X = (x _ {i j}), x = (x _ {i}))\tag{5}
\]

是 G 上的左 Haar 测度。另一方面，根据 §2, No.~9 的 Prop. 14，

\[
\Delta_ {\mathrm{G}} ((X, x)) = \Delta_ {\mathbf {G L} (n, K)} (X) \Delta_ {K ^ {n}} (x) (\mathrm{mod} \det X) ^ {- 1},
\]

或

\[
\Delta_ {\mathrm{G}} ((X, x)) = \operatorname{mod} (\det X ^ {- 1}).\tag{6}
\]

因此，G 上的一个右 Haar 测度为

\[
(\mathrm{mod} \det X) ^ {- n} \cdot \left(\bigotimes_ {i j} d x _ {i j}\right) \otimes \left(\bigotimes_ {i} d x _ {i}\right).\tag{7}
\]

例 3. --- 严格三角群。

令 \([1, n]\) 为整数 \(m\) 的集合，其中满足 \(1 \leqslant m \leqslant n\)。令 \(J\) 为 \([1, n] \times [1, n]\) 的一个子集，满足以下条件：

\begin{enumerate}
\def\labelenumi{\arabic{enumi})}
\item
  若 \((i,j)\in \mathbf{J}\)，则 \(i <   j\)
\item
  若 \((i,j)\notin J\)，则对每个整数 \(k\) 满足 \(i < k < j\)，\((i,k)\) 和 \((k,j)\) 这两个有序对中至少有一个不属于 \(J\)。
\end{enumerate}

令 \(T_J\) 为如下矩阵的集合：\(Z = (z_{ij})_{1 \leqslant i \leqslant n, 1 \leqslant j \leqslant n}\) 的元素属于 \(K\)，且 \(z_{ii} = 1\)，并且 \(z_{ij} = 0\)，当 \(i \neq j\) 且 \((i,j) \notin J\) 时。这是 \(\mathbf{GL}(n,K)\) 的闭子集。映射 \(Z \mapsto (z_{ij})_{(i,j) \in J}\) 是从 \(T_J\) 到 \(K^s\) 的同胚（其中 \(s\) 表示 \(J\) 的元素个数）。若 \(Z' = (z_{ij}') \in T_J\)，则 \(Z'Z = (z_{ij}'')\)，其中

\[
\begin{array}{c} z _ {i j} ^ {\prime \prime} = z _ {i j} + z _ {i j} ^ {\prime} + \sum_ {i <   h <   j} z _ {i h} ^ {\prime} z _ {h j} \quad \text {for} i <   j, \\ z _ {i j} ^ {\prime \prime} = 0 \text {for} i > j, \quad z _ {i i} ^ {\prime \prime} = 1, \end{array}
\]

从而 \(Z^{\prime}Z \in T_{J}\)。若将 \(T_{J}\) 与 \(K^{s}\) 认同，则映射 \(Z \mapsto Z^{\prime}Z\)（\(Z^{\prime}\) 固定时）对应于一个仿射映射，其行列式为 1；将 \((i,j) \in J\) 中的有序对按字典序排列并应用以下引理即可看出这一点：

引理 6. --- 令 L 为一个全序有限集。对每个 \(\lambda \in L\)，令 \(V_{\lambda}\) 为交换环 \(k\) 上的有限维自由模；对于 L 中的 \(\lambda, \mu\)，若满足 \(\lambda \leqslant \mu\)，令 \(f_{\lambda \mu} \in \operatorname{Hom}_k(V_\mu, V_\lambda)\)。则线性映射

\[
(v _ {\lambda}) _ {\lambda \in \mathrm{L}} \mapsto \left(\sum_ {\mu \geqslant \lambda} f _ {\lambda \mu} (v _ {\mu})\right) _ {\lambda \in \mathrm{L}},
\]

从 \(\prod_{\lambda \in L}V_{\lambda}\) 到 \(\prod_{\lambda \in L}V_{\lambda}\)，其行列式为 \(\prod_{\lambda \in L}\det f_{\lambda \lambda}\)。

立即可化归到 L 是整数区间的情形，此时引理由 A, III, §8, No.~6 的公式 (31) 得出。

若 \(Z \in \mathrm{T}_{\mathrm{J}}\)，则可见存在 \(Z' \in \mathrm{T}_{\mathrm{J}}\) 使得 \(Z'Z = I_n\)，从而 \(Z' = Z^{-1}\)。因此，\(\mathrm{T}_{\mathrm{J}}\) 是 \(\mathbf{GL}(n,\mathrm{K})\) 的闭子群。另一方面，§1, No.~10 的 Prop. 15 表明，测度

\[
\bigotimes_ {(i, j) \in J} d z _ {i j}
\]

是 \(T_{J}\) 上的左 Haar 测度。通过计算 \(ZZ'\)，同理可见此测度也是 \(T_{J}\) 上的右 Haar 测度。

如果在 \(T_{J}\) 的定义中交换行与列的作用，则有类似的结果。

当 J 是有序对 \((i,j)\) 的集合，且满足 i\textless j 时，群 \(T_{J}\) 称为 K 上 n 阶上严格三角群，记为 \(\mathrm{T}_{1}(n,\mathrm{K})\)。其转置称为下严格三角群。

例 4. --- 大三角群。

令 \(n_1, \ldots, n_r\) 为满足 \(\geqslant 1\) 的整数。置 \(p_k = n_1 + \ldots + n_{k-1}\)，\(n = p_{r+1} = n_1 + \cdots + n_r\)。令 \(I_k\) 为整数 \(j\) 的集合，其中满足 \(p_k < j \leqslant p_{k+1}\)，并令 \(J\) 为所有 \(I_k \times I_l\)（其中 \(k < l\)）的并。令 \(G\) 为 \(\mathbf{GL}(n, K)\) 的闭子群，其元素是满足下列条件的矩阵 \((Z_{kl})_{1 \leqslant k \leqslant r, 1 \leqslant l \leqslant r}\)：

\begin{enumerate}
\def\labelenumi{\arabic{enumi})}
\item
  每个 \(Z_{kl}\) 是矩阵 \((z_{ij})_{i\in \mathrm{I}_k,j\in \mathrm{I}_l}\)，其元素属于 \(\mathbf{K}\)，有 \(n_k\) 行和 \(n_l\) 列；
\item
  \(Z_{kl} = 0\)，当 \(k > l\) 时；
\item
  \(Z_{kk} \in \mathbf{GL}(n_k, \mathrm{K})\)，其中 \(1 \leqslant k \leqslant r\)。
\end{enumerate}

分块乘法公式

\[
\begin{array}{r} \left( \begin{array}{c c c c} Z _ {1 1} & 0 & \ldots & 0 \\ 0 & Z _ {2 2} & \ldots & 0 \\ \vdots & \vdots & \ddots & \vdots \\ 0 & 0 & \ldots & Z _ {r r} \end{array} \right) \left( \begin{array}{c c c c} 1 & Z _ {1 2} & \ldots & Z _ {1 r} \\ 0 & 1 & \ldots & Z _ {2 r} \\ \vdots & \vdots & \ddots & \vdots \\ 0 & 0 & \ldots & 1 \end{array} \right) \\ = \left( \begin{array}{c c c c} Z _ {1 1} & Z _ {1 1} Z _ {1 2} & \ldots & Z _ {1 1} Z _ {1 r} \\ 0 & Z _ {2 2} & \ldots & Z _ {2 2} Z _ {2 r} \\ \vdots & \vdots & \ddots & \vdots \\ 0 & 0 & \ldots & Z _ {r r} \end{array} \right) \end{array}\tag{8}
\]

表明 G 是由元素 \((Z_{kl}) \in G\)（满足 \(Z_{kl} = 0\)，其中 \(k \neq l\)）所成的子群 D 与例 3 的子群 \(T_{J}\) 的拓扑半直积。此外，D 同构于各群 \(\mathbf{GL}(n_{k}, \mathbf{K})\)（\(1 \leqslant k \leqslant r\)）的直积。

令 \(J'\) 为有序对 \((j, i)\) 的集合，其中 \((i, j) \in J\)，并令 H 为 \((i, j) \in [1, n] \times [1, n]\) 中不属于 \(J'\) 的有序对的集合。令 \(Z' = (z_{ij})_{1 \leqslant i \leqslant n, 1 \leqslant j \leqslant n}\) 为 G 的一个元素。根据 §2, No.~9 的 Prop. 14 以及上面的例 1 和例 3，取以下测度在映射下的像即可得到 G 上的一个左 Haar 测度：

\[
\bigotimes_ {k = 1} ^ {r} \left(\left(\operatorname{mod} \det Z _ {k k}\right) ^ {- n _ {k}} \cdot \bigotimes_ {i, j \in I _ {k}} d z _ {i j}\right) \otimes \left(\bigotimes_ {(i, j) \in J} d z _ {i j}\right)
\]

在映射

\[
\bigl ((Z _ {k k}), (Z _ {k l}) \bigr) \mapsto \left( \begin{array}{c c c c} Z _ {1 1} & Z _ {1 1} Z _ {1 2} & \ldots & Z _ {1 1} Z _ {1 r} \\ 0 & Z _ {2 2} & \ldots & Z _ {2 2} Z _ {2 r} \\ \vdots & \vdots & \ddots & \vdots \\ 0 & 0 & \ldots & Z _ {r r} \end{array} \right).
\]

现在，对 \(k < l\)，考虑矩阵 \(Z_{kl} = (z_{ij})_{i \in \mathbf{I}_k, j \in \mathbf{I}_l}\) 所成的向量空间。它是 \(n_l\) 个子空间 \(\mathbf{M}_j\)（\(j \in \mathbf{I}_l\)）的直和；这些子空间由满足 \(z_{ih} = 0\)（\(h \neq j\)）的矩阵组成。每个子空间 \(\mathbf{M}_j\) 在映射 \(Z_{kl} \mapsto Z_{kk}Z_{kl}\) 下稳定，并且该映射在 \(\mathbf{M}_j\) 上的限制具有矩阵 \(Z_{kk}\)。因此（§1, No.~10, Prop. 15），测度

\(\bigotimes_{i \in \mathbf{I}_k, j \in \mathbf{I}_l} dz_{ij}\) 在映射 \(Z_{kl} \mapsto Z_{kk} Z_{kl}\) 下的像为

\[
(\mathrm{mod} \det Z _ {k k}) ^ {- n _ {l}} \cdot \bigotimes_ {i \in \mathrm{I} _ {k}, j \in \mathrm{I} _ {l}} d z _ {i j}.
\]

因此，G 上的一个左 Haar 测度为

\[
\prod_ {k = 1} ^ {r} (\mathrm{mod} \det Z _ {k k}) ^ {- q _ {k}} \cdot \bigotimes_ {(i, j) \in \mathrm{H}} d z _ {i j}\tag{9}
\]

其中 \(q_{k} = \sum_{k\leqslant l\leqslant r}n_{l} = n - p_{k}\)

我们再次利用 §2 的 Prop. 14 来计算 G 的模。群 D 和 T \(_{J}\) 是幺模的；另一方面：

\[
\begin{array}{c} \left( \begin{array}{c c c c} Z _ {1 1} & 0 & \ldots & 0 \\ 0 & Z _ {2 2} & \ldots & 0 \\ \vdots & \vdots & \ddots & \vdots \\ 0 & 0 & \ldots & Z _ {r r} \end{array} \right) \left( \begin{array}{c c c c} 1 & Z _ {1 2} & \ldots & Z _ {1 r} \\ 0 & 1 & \ldots & Z _ {2 r} \\ \vdots & \vdots & \ddots & \vdots \\ 0 & 0 & \ldots & 1 \end{array} \right) \left( \begin{array}{c c c c} Z _ {1 1} & 0 & \ldots & 0 \\ 0 & Z _ {2 2} & \ldots & 0 \\ \vdots & \vdots & \ddots & \vdots \\ 0 & 0 & \ldots & Z _ {r r} \end{array} \right) ^ {- 1} \\ = \left( \begin{array}{c c c c} 1 & Z _ {1 2} ^ {\prime} & \ldots & Z _ {1 r} ^ {\prime} \\ 0 & 1 & \ldots & Z _ {2 r} ^ {\prime} \\ \vdots & \vdots & \ddots & \vdots \\ 0 & 0 & \ldots & 1 \end{array} \right), \end{array}
\]

其中 \(Z_{kl}^{\prime}=Z_{kk}Z_{kl}Z_{ll}^{-1}\)。考虑例 3 以及 §1, No.~10 的 Prop. 15，并像上面那样论证可见，若 \(X=\mathrm{diag}(Z_{11},\ldots,Z_{rr})\in\mathbf{D}\)，则自同构 \(Z\mapsto X^{-1}ZX\) 的模（在 \(T_{J}\) 上）为

\[
\prod_ {k <   l} (\mathrm{mod} \det Z _ {k k}) ^ {- n _ {l}} (\mathrm{mod} \det Z _ {l l}) ^ {n _ {k}},
\]

因此

\[
\Delta_ {\mathrm{G}} (Z) = \prod_ {k = 1} ^ {r} (\mathrm{mod} \det Z _ {k k}) ^ {n + n _ {k} - 2 q _ {k}}.\tag{10}
\]

G 的转置群 \(G'\) 以同样方式研究。对于 \(G'\)，可得左 Haar 测度

\[
\prod_ {k = 1} ^ {r} (\mathrm{mod} \det Z _ {k k}) ^ {- p _ {k + 1}} \cdot \bigotimes_ {(j, i) \in \mathrm{H}} d z _ {i j},
\]

以及模

\[
\prod_ {k = 1} ^ {r} (\mathrm{mod} \det Z _ {k k}) ^ {n + n _ {k} - 2 p _ {k + 1}}.
\]

特别地，若取 \(n_{1}=\ldots=n_{r}=1\)，则所得的群 G 是 \(T(n,K)^{*}\)——即子代数 \(\mathbf{M}_{n}(\mathbf{K})\) 的可逆元群，该子代数由矩阵 \(X=(x_{ij})\) 构成，并满足 \(x_{ij}=0\) 当 i\textgreater j 时。这个代数记为 \(\mathrm{T}(n,\mathrm{K})\)，称为上三角代数，而群 \(\mathrm{T}(n,\mathrm{K})^{*}\) 称为 K 上 n 阶上大三角群。于是前面的公式成为：\(\mathrm{T}(n,\mathrm{K})^{*}\) 上的一个左 Haar 测度为

\[
\prod_ {i = 1} ^ {n} (\mathrm{mod} z _ {i i}) ^ {i - n - 1} \cdot \bigotimes_ {i \leqslant j} d z _ {i j} \quad (Z = (z _ {i j}))\tag{9 bis}
\]

而 \(\mathrm{T}(n,\mathbf{K})^*\) 的模为

(10 bis)

\[
\Delta_ {\mathrm{T} (n, \mathrm{K}) ^ {*}} (Z) = \prod_ {i = 1} ^ {n} (\mathrm{mod} z _ {i i}) ^ {2 i - n - 1} \quad \left(Z = \left(z _ {i j}\right)\right).
\]

对于 \(\mathrm{T}(n,\mathbf{K})^*\) 的转置群，即下大三角群，可得左 Haar 测度

\[
\prod_ {i = 1} ^ {n} (\mathrm{mod} z _ {i i}) ^ {- i} \cdot \bigotimes_ {i \geqslant j} d z _ {i j},
\]

以及模

\[
\prod_ {i = 1} ^ {n} (\mathrm{mod} z _ {i i}) ^ {n + 1 - 2 i}.
\]

备注。--- 群 \(\mathrm{T}(n,\mathrm{K})^{*}\) 是 \(\mathbf{GL}(n,\mathrm{K})\) 的闭子群，且 \(\Delta_{\mathrm{T}(n,\mathrm{K})^{*}}\left((z_{ij})\right)=\prod_{i=1}^{n}(\mathrm{mod}z_{ii})^{2i-n-1}\)。我们在例 1 中看到 \(\Delta_{\mathbf{GL}(n,\mathrm{K})}=1\)。若 n\textgreater1，函数

\[
\Delta_ {\mathrm{T} (n, \mathrm{K}) ^ {*}} / \Delta_ {\mathbf {G L} (n, \mathrm{K})}
\]

在 \(\mathrm{T}(n,\mathbf{K})^*\) 上不能延拓为 \(\mathbf{GL}(n,\mathbf{K})\) 到 \(\mathbf{C}^*\) 的连续表示（因为根据引理 5，这样的表示在 \(\mathbf{SL}(n,\mathbf{K})\) 上应等于 1，而 \(\mathrm{mod}(z_{11})^{1 - n}\neq 1\) 对适当选取的 \(z_{11}\) 成立）。由此可知，齐性空间 \(\mathbf{GL}(n,\mathbf{K}) / \mathrm{T}(n,\mathbf{K})^*\) 当 \(n > 1\) 时不存在相对不变测度（\(\S 2\)，No.~6, Th. 3 的 Cor. 1）。

当 n = 2 时，这个齐性空间可认同为 K 上的射影直线。事实上，令 \((e_{1}, e_{2})\) 为 \(K^{2}\) 的典范基。群 \(\mathbf{GL}(2, \mathbf{K})\) 在去掉 0 的 \(K^{2}\) 的直线集上传递地作用，而 \(Ke_{1} - \{0\}\) 的稳定子群是 \(T(2, K)^{*}\)。

例 5. --- 特殊三角群。

我们重新采用例 4 开头的记号，并考虑子群 \(\mathbf{G}_{1} = \mathbf{G} \cap \mathbf{SL}(n, \mathbf{K})\)。这个子群是群 \(\mathbf{D}_{1} = \mathbf{D} \cap \mathbf{SL}(n, \mathbf{K})\) 与 \(T_{J}\) 的拓扑半直积。群 \(D_{1}\) 有一个同构于 \(\mathbf{SL}(n_{r}, \mathbf{K})\) 的正规子群 A，即由元素 \(\text{diag}(Z_{kk})\)（其中 \(Z_{kk} = 1\) 对 k \textless{} r）组成的子群。下面的同态

\[
\varphi : \operatorname{diag} (Z _ {1 1}, \dots , Z _ {r r}) \mapsto (Z _ {1 1}, \dots , Z _ {r - 1, r - 1})
\]

从 \(D_{1}\) 到 \(\mathbf{GL}(n_{1},\mathbf{K})\times\cdots\times\mathbf{GL}(n_{r-1},\mathbf{K})\) 是满射且核为 A。另一方面，\(\varphi\) 是连续的。考虑附录 I 的引理 2，可将 \(D_{1}/A\) 与 \(\mathbf{GL}(n_{1},\mathbf{K})\times\cdots\times\mathbf{GL}(n_{r-1},\mathbf{K})\) 认同。我们将 A 的 Haar 测度记为 \(\mu\)（cf.~例 6），并将

\[
\alpha = \bigotimes_ {k = 1} ^ {r - 1} \left(\left(\operatorname{mod} \det Z _ {k k}\right) ^ {- n _ {k}} \cdot \bigotimes_ {i, j \in I _ {k}} d z _ {i j}\right) \otimes^ {\prime} d \mu \left(Z _ {r r}\right)
\]

满足以下条件的 \(\mathbf{D}_1\) 上的 Haar 测度

\[
\alpha / \mu = \bigotimes_ {k = 1} ^ {r - 1} \left(\left(\operatorname{mod} \det Z _ {k k}\right) ^ {- n _ {k}} \cdot \bigotimes_ {i, j \in I _ {k}} d z _ {i j}\right)
\]

（§2, No.~7, Prop. 10）。然后如例 4 所示，可得 \(\mathbf{G}_1\) 上的一个左 Haar 测度为

\[
\begin{array}{l} \bmod \Big (\prod_ {k = 1} ^ {r - 1} (\det Z _ {k k}) ^ {n _ {k} - q _ {k}} \Big) \\ \qquad \cdot \left[ \bigotimes_ {k = 1} ^ {r - 1} \big ((\bmod \det Z _ {k k}) ^ {- n _ {k}} \cdot \bigotimes_ {i, j \in I _ {k}} d z _ {i j} \big) \otimes^ {\prime} d \mu (Z _ {r r}) \right] \otimes \bigotimes_ {(i, j) \in J} d z _ {i j}. \end{array}
\]

由于 \(\mathbf{G}_1\) 是 \(\mathbf{G}\) 的正规子群，\(\mathbf{G}_1\) 的模是 \(\mathbf{G}\) 的模的限制（§2, No.~7, Prop. 10 b)）。

若 \(n_{r}=1\)，子群 A 退化为单位元，G 上的一个左 Haar 测度为

\[
\operatorname{mod} \left(\prod_ {k = 1} ^ {r - 1} (\det Z _ {k k}) ^ {- q _ {k}}\right) \cdot \bigotimes_ {k = 1} ^ {r - 1} \left(\bigotimes_ {i, j \in I _ {k}} d z _ {i j}\right) \otimes \bigotimes_ {(i, j) \in J} d z _ {i j}.
\]

若取 \(n_{1}=n_{2}=\ldots=n_{r}=1\)，所得群 \(G_{1}\) 称为上特殊三角群，其转置 \(G_{1}^{\prime}\) 称为下特殊三角群。\(G_{1}\) 上的一个左 Haar 测度为

\[
\operatorname{mod} \left(\prod_ {i = 1} ^ {n - 1} z _ {i i} ^ {i - n - 1}\right) \cdot \left(\bigotimes_ {i = 1} ^ {n - 1} d z _ {i i}\right) \otimes \left(\bigotimes_ {1 \leqslant i <   j \leqslant n} d z _ {i j}\right)\tag{11}
\]

而 \(\mathbf{G}_1\) 的模为

\[
\operatorname{mod} \left(\prod_ {i = 1} ^ {n - 1} z _ {i i} ^ {2 i - 2 n}\right).\tag{12}
\]

对于 \(G_{1}^{\prime}\)，同样可得左 Haar 测度

\[
\operatorname{mod} \left(\prod_ {i = 1} ^ {n - 1} z _ {i i} ^ {n - i - 1}\right) \cdot \left(\bigotimes_ {i = 1} ^ {n - 1} d z _ {i i}\right) \otimes \left(\bigotimes_ {1 \leqslant j <   i \leqslant n} d z _ {i j}\right)
\]

以及模

\[
\operatorname{mod} \biggl (\prod_ {i = 1} ^ {n - 1} z _ {i i} ^ {2 n - 2 i} \biggr).
\]

例 6. --- 特殊线性群。

\(\mathrm{T}_{1}(n,\mathrm{K})\) 和 \({}^{t}(\mathrm{T}(n,\mathrm{K})^{*})\) 是 \(\mathbf{GL}(n,\mathbf{K})\) 的闭子群，且二者的交为 \(\{e\}\)。因此，映射 \((M,N)\mapsto M\cdot N\) 是一个连续双射 \(\varphi\)，其定义域为 \(\mathrm{T}_{1}(n,\mathrm{K})\times{}^{t}(\mathrm{T}(n,\mathrm{K})^{*})\)，值域为子集 \(\Omega\) 的 \(\mathbf{GL}(n,\mathbf{K})\)。

引理 7. --- a) 令 \(U = (u_{ij}) \in \mathbf{GL}(n, K)\)。要使 \(U \in \Omega\)，充分必要条件是 \(\det(u_{ij})_{k \leqslant i,j \leqslant n} \neq 0\) 对 \(k = 2, 3, \ldots, n\) 成立。

\begin{enumerate}
\def\labelenumi{\alph{enumi})}
\setcounter{enumi}{1}
\item
  \(\Omega\) 是 \(\mathbf{GL}(n,\mathbf{K})\) 的开子集。
\item
  映射 \(\varphi\) 是从 \(\mathrm{T}_1(n,\mathrm{K})\times {}^t (\mathrm{T}(n,\mathrm{K})^*)\) 到 \(\Omega\) 的同胚。
\end{enumerate}

要使 \(U \in \Omega\)，充分必要条件是存在 \(Z = (z_{ij}) \in \mathrm{T}_1(n,\mathrm{K})\)，使得 \(ZU \in {}^t (\mathbf{T}(n,\mathbf{K}))\)（此时必然

\(ZU \in {}^{t}(\mathrm{T}(n,\mathrm{K})^{*})\)，因为 \(U\) 和 \(Z\) 都可逆）。由前面的结论，若 \(Z\) 存在，则它是唯一的。因此，要使 \(U \in \Omega\)，充分必要条件是线性方程组

\[
\sum_ {k = 1} ^ {n} z _ {i k} u _ {k j} = 0 \quad (1 \leqslant i <   j \leqslant n)
\]

（其中 \((z_{ij})\in \mathrm{T}_1(n,\mathbf{K})\)）有唯一解。这个方程组可以写成

\[
\sum_ {k = i + 1} ^ {n} z _ {i k} u _ {k j} = - u _ {i j} \quad (1 \leqslant i <   j \leqslant n).\tag{13}
\]

固定 \(i\)，关于 \(n - i\) 个未知数 \(z_{i,i+1}\)、\(z_{i,i+2}\)、\ldots、\(z_{i,n}\) 有一个方程组；这些方程组有唯一解的充分必要条件是

\[
\det (u _ {k j}) _ {i + 1 \leqslant k \leqslant n, i + 1 \leqslant j \leqslant n} \neq 0
\]

对 \(i = 1,2,\ldots ,n - 1\) 成立。这证明了 a)。由此可知 \(\Omega\) 在 \(\mathbf{GL}(n,\mathbf{K})\) 中是开集。另一方面，利用克拉默公式求解方程组 (13) 时，\(z_{ij}\) 是 \(u_{ij}\) 的有理函数，且在 \(\Omega\) 中分母非零，因此 \(Z\) 连续依赖于 \(U\)（在 \(\Omega\) 中），这证明了 c)。

现在令 \(\mathbf{G}_{1}^{\prime}\subset{}^{t}(\mathrm{T}(n,\mathrm{K})^{*})\) 为下特殊三角群。映射 \((M,N)\mapsto M\cdot N\) 是连续双射 \(\psi\)，其定义域为 \(\mathrm{T}_{1}(n,\mathrm{K})\times\mathbf{G}_{1}^{\prime}\)，值域为子集 \(\Omega^{\prime}\) 的 \(\mathbf{SL}(n,\mathbf{K})\)。

引理 8. --- a) 令 \(U = (u_{ij}) \in \mathbf{SL}(n, \mathbf{K})\)。要使 \(U \in \Omega'\)，充分必要条件是 \(\det(u_{ij})_{k \leqslant i,j \leqslant n} \neq 0\) 对 \(k = 2,3,\ldots,n\) 成立。

\begin{enumerate}
\def\labelenumi{\alph{enumi})}
\setcounter{enumi}{1}
\item
  \(\Omega'\) 是 \(\mathbf{SL}(n,\mathbf{K})\) 的开子集。
\item
  映射 \(\psi\) 是从 \(\mathrm{T}_1(n,\mathbf{K})\times \mathrm{G}_1'\) 到 \(\Omega^{\prime}\) 的同胚。事实上，令 \(M\in \mathrm{T}_1(n,\mathbf{K})\) 且 \(N\in {}^t (\mathrm{T}(n,\mathbf{K})^*)\)。要使 \(M\cdot N\in \mathbf{SL}(n,\mathbf{K})\)，充分必要条件是 \(N\in \mathbf{G}_1'\)。因此 \(\Omega^{\prime} = \mathbf{SL}(n,\mathbf{K})\cap \Omega\)，引理 8 立即由引理 7 得出。
\end{enumerate}

命题 6. --- a) 群 \(\mathbf{SL}(n,\mathbf{K})\) 是幺模的。

\begin{enumerate}
\def\labelenumi{\alph{enumi})}
\setcounter{enumi}{1}
\item
  令 \(\mu_1\) 和 \(\mu_2\) 分别为上严格三角群 \(\mathrm{T}_1(n, \mathrm{K})\) 和下特殊三角群 \(\mathrm{G}_1'\) 上的左 Haar 测度。\(\mu_1 \otimes \mu_2\) 在同胚 \((M, N) \mapsto M \cdot N^{-1}\) 下的像（该同胚从 \(\mathrm{T}_1(n, \mathrm{K}) \times \mathrm{G}_1'\) 到 \(\Omega'\)）是 \(\Omega'\) 上某个 \(\mathrm{SL}(n, \mathrm{K})\) Haar 测度的限制。
\item
  \(\Omega'\) 在 \(\mathbf{SL}(n,\mathbf{K})\) 中的补集对于 \(\mathbf{SL}(n,\mathbf{K})\) 的 Haar 测度是零测的。
\end{enumerate}

群 \(\mathbf{GL}(n,\mathbf{K})\) 是幺模的（例 1），而 \(\mathbf{SL}(n,\mathbf{K})\) 是 \(\mathbf{GL}(n,\mathbf{K})\) 的正规子群，因此也是幺模的（ \(\S 2\) ，No.~7, Prop. 10 b)）。断言 b) 由 a)、引理 8 以及 \(\S 2\) ，No.~9 的 Prop. 13 得出。我们来证明 c)。根据引理 8 a)，只需证明如下事实：若 \(p((u_{ij})_{1 \leqslant i,j \leqslant n})\) 是一个在 \(\mathbf{SL}(n,\mathbf{K})\) 上不恒为零的多项式，则满足 \(U \in \mathbf{SL}(n,\mathbf{K})\) 且 \(p(U) = 0\) 的元素所成的集合 E 对 Haar 测度是零测的。考虑 \(\S 1\) ，No.~10, Cor. of Prop. 13，\(\mathbf{SL}(n,\mathbf{K})\) 的拓扑有一个可数基。因此，只需证明：对每个 \(U_0 \in \mathbf{E}\)，存在 \(U_0\) 在 \(\mathbf{SL}(n,\mathbf{K})\) 中的一个邻域，其与 E 的交是零测的；或者说，存在 \(\mathbf{SL}(n,\mathbf{K})\) 中 I 的一个邻域 W，使得 \(U_0^{-1}\mathbf{E} \cap \mathbf{W}\) 是零测的。令 \(\mathbf{W} = \Omega'\)。由 b) 可知，问题归结为证明：对 \((M,N) \in \mathrm{T}_1(n,\mathbf{K}) \times \mathrm{G}_1'\)，满足 \(p(U_0MN) = 0\) 的元素所成的集合对于 \(\mu_1 \otimes \mu_2\) 是零测的。根据 \(\mu_1\) 和 \(\mu_2\) 的表达式（在例 3 和例 5 中计算），这将由以下引理得到：

引理 9. --- 令 \(\psi\) 为 \(\neq 0\) 的 \(\mathbf{K}[\mathbf{X}_1, \ldots, \mathbf{X}_r]\) 中的一个多项式。在空间 \(\mathbf{K}^r\) 中，集合 \(\mathbf{N}\) 由 \(\psi(x_1, \ldots, x_r) = 0\) 定义，对 Haar 测度是零测的。

我们对 r 作归纳论证。当 r = 1 时，引理显然成立，因为此时 N 是有限集。如有必要改变变量的编号，可设 \(\psi \notin K[X_{1}, \ldots, X_{r-1}]\)；写成

\[
\psi \left(\mathrm{X} _ {1}, \dots , \mathrm{X} _ {r}\right) = \mathrm{X} _ {r} ^ {m} \psi_ {0} \left(\mathrm{X} _ {1}, \dots , \mathrm{X} _ {r - 1}\right) + \dots + \psi_ {m} \left(\mathrm{X} _ {1}, \dots , \mathrm{X} _ {r - 1}\right)
\]

其中 \(m > 0\) 且 \(\psi_0 \neq 0\)。在空间 \(\mathbf{K}^{r-1}\) 中，令 \(\mathbf{N}_0\) 为由 \(\psi_0(x_1, \ldots, x_{r-1}) = 0\) 定义的集合，根据归纳假设它是零测的。对于每个 \((x_1, \ldots, x_{r-1}) \notin \mathbf{N}_0\)，\(x_r \in \mathbf{K}\) 中使得 \((x_1, \ldots, x_{r-1}, x_r) \in \mathbf{N}\) 的集合是有限的，因而是零测的。由于 \(\mathbf{K}^r\) 在无穷远处可数（§1, No.~10, Prop. 13 的 Cor.），\(\mathbf{N} \cap [(\mathbf{K}^{r-1} - \mathbf{N}_0) \times \mathbf{K}]\) 在 \(\mathbf{K}^r\) 中是零测的。因此 \(\mathbf{N}\) 是零测的。

例 7. --- Iwasawa 分解 \(\mathbf{GL}(n,\mathbf{K})\).

在这个例子中，K 总表示域 R、C、H 中的一个。若 \(\lambda \in K\)，则规定：当 K = R 时 \(\overline{\lambda}\) 等于 \(\lambda\)，当 K = C 或 H 时取 \(\lambda\) 的共轭。令 E 为 K 上的右向量空间，维数为 n，并令 \(\Phi\) 为 E 上的非退化正厄米型。

引理 10. --- 设 \((f_1, f_2, \ldots, f_n)\) 为 E 的一个基。

\begin{enumerate}
\def\labelenumi{\alph{enumi})}
\item
  存在唯一的 E 的标准正交基 \((e_{1}, e_{2}, \ldots, e_{n})\)，使得 \(f_{i} = e_{1}\alpha_{i1} + e_{2}\alpha_{i2} + \cdots + e_{i}\alpha_{ii}\)（\(i = 1, 2, \ldots, n\)），且对所有 i 都有 \(\alpha_{ii} > 0\)。
\item
  对于固定的 \(\Phi\)，\(e_i\) 和 \(\alpha_{ij}\) 连续依赖于 \((f_1, \ldots, f_n) \in \mathbf{E}^n\)。
\end{enumerate}

令 \(E_{i}=f_{1}K+f_{2}K+\cdots+f_{i}K\)，它的维数为 i。令 \(g_{i}\) 为 \(E_{i}\) 中一个非零元素，正交于 \(E_{i-1}\)，且满足 \(\Phi(g_{i},g_{i})=1\)。对 i 作归纳可见，\((g_{1},\ldots,g_{i})\) 是 \(E_{i}\) 的一个标准正交基。特别地，\((g_{1},\ldots,g_{n})\) 是 E 的一个标准正交基。令 \(\lambda_{i}=\Phi(f_{i},g_{i})\)。由于 \(f_{i}\notin E_{i-1}\)，有 \(\lambda_{i}\neq0\)。置 \(e_{i}=g_{i}\left|\lambda_{i}\right|\lambda_{i}^{-1}\)。则

\[
\Phi (e _ {i}, e _ {i}) = | \lambda_ {i} | ^ {2} \overline {{{{\lambda}}}} _ {i} ^ {- 1} \Phi (g _ {i}, g _ {i}) \lambda_ {i} ^ {- 1} = 1,
\]

因此 \((e_1, \ldots, e_i)\) 也是 \(\mathbf{E}_i\) 的一个标准正交基；此外，\(\Phi(e_i, f_i) = |\lambda_i| \overline{\lambda}_i^{-1} \Phi(g_i, f_i) = |\lambda_i| > 0\)，所以 \(e_i\) 具有 a) 中的性质。令 \((e_1', \ldots, e_n')\) 为具有相同性质的 \(\mathbf{E}\) 的另一个标准正交基。由 \(i\) 作归纳可见，\((e_1', \ldots, e_i')\) 必为 \(\mathbf{E}_i\) 的一个基，因此 \(e_i' = e_i \mu_i\)，其中 \(\mu_i \in \mathbf{K}\)。于是

\[
1 = \Phi (e _ {i} ^ {\prime}, e _ {i} ^ {\prime}) = \overline {{\mu}} _ {i} \Phi (e _ {i}, e _ {i}) \mu_ {i} = \overline {{\mu}} _ {i} \mu_ {i},
\]

且 \(0 < \Phi(e_i', f_i) = \overline{\mu}_i \Phi(e_i, f_i)\)，所以 \(\mu_i > 0\) 且 \(\mu_i^2 = 1\)，从而 \(\mu_i = 1\)，这就证明了 a)。假定已经证明 \(e_i\) 和 \(\alpha_{ij}\) 连续依赖于 \((f_1, \ldots, f_n)\)，其中 \(i < i_0\)，现证明 \(e_{i_0}\) 和 \(\alpha_{i_0j}\) 连续依赖于 \((f_1, \ldots, f_n)\)。当 \(j < i_0\) 时，\(\overline{\alpha}_{i_0j} = \Phi(f_{i_0}, e_j)\) 根据归纳假设连续依赖于 \((f_1, \ldots, f_n)\)。另一方面，

\[
\Phi (f _ {i _ {0}}, f _ {i _ {0}}) = | \alpha_ {i _ {0} 1} | ^ {2} + | \alpha_ {i _ {0} 2} | ^ {2} + \dots + | \alpha_ {i _ {0}, i _ {0} - 1} | ^ {2} + \alpha_ {i _ {0} i _ {0}} ^ {2},
\]

因此 \(\alpha_{i_0i_0}\) 连续依赖于 \((f_1,\ldots ,f_n)\)。因此

\[
e _ {i _ {0}} = (f _ {i _ {0}} - e _ {1} \alpha_ {i _ {0} 1} - \dots - e _ {i _ {0} - 1} \alpha_ {i _ {0}, i _ {0} - 1}) \alpha_ {i _ {0} i _ {0}} ^ {- 1}
\]

连续依赖于 \((f_1, \ldots, f_n)\)。

以下设 \(\mathbf{E} = \mathbf{K}^n\)，并取 \(\Phi\) 为如下形式

\[
\overline {{{x}}} _ {1} y _ {1} + \dots + \overline {{{x}}} _ {n} y _ {n}.
\]

回忆一下，\(\mathbf{U}(n,\mathbf{K})\) 表示相应的酉群。即使 K 不可交换，我们仍将 \(\mathrm{T}_{1}(n,\mathrm{K})\) 记为 \(\mathbf{M}_{n}(\mathrm{K})\) 的上三角矩阵群，其中所有对角元素均为 1。

命题 7. --- 令 \(D_{+}^{*}\) 为对角矩阵的群，其对角元素均 \textgreater0。映射 \((U,D,T)\mapsto UDT\) 是从 \(\mathbf{U}(n,\mathbf{K})\times\mathbf{D}_{+}^{*}\times\mathbf{T}_{1}(n,\mathbf{K})\) 到 \(\mathbf{GL}(n,\mathbf{K})\) 的同胚。

令 \((\varepsilon_1, \ldots, \varepsilon_n)\) 为 \(\mathbf{K}^n\) 的典范基。令 \(X \in \mathbf{GL}(n, \mathbf{K})\)。则 \(X \cdot \varepsilon_i = f_i\) 构成 \(\mathbf{E}\) 的一个基。根据引理 10，可以将这个基 \((f_i)\) 与一个基 \((e_i)\) 联系起来。令 \(U\) 为 \(\mathbf{E}\) 的酉自同构的矩阵，该自同构将 \(\varepsilon_i\) 变为 \(e_i\)。于是

\[
U ^ {- 1} \cdot f _ {i} = \varepsilon_ {1} \alpha_ {i 1} + \varepsilon_ {2} \alpha_ {i 2} + \dots + \varepsilon_ {i} \alpha_ {i i}
\]

其中 \(\alpha_{ii} > 0\) 对 \(i = 1,2,\ldots ,n\) 成立。于是 \(X = UC\)，其中 \(C\) 是矩阵

\[
\left( \begin{array}{c c c c} \alpha_ {1 1} & \alpha_ {2 1} & \ldots & \alpha_ {n 1} \\ 0 & \alpha_ {2 2} & \ldots & \alpha_ {n 2} \\ \vdots & \vdots & \ddots & \vdots \\ 0 & 0 & \ldots & \alpha_ {n n} \end{array} \right).
\]

此外，根据引理 10，U 和 C 连续依赖于 X。另一方面，公式 (8) 表明 C 可写成 DT 的形式，其中 \(D \in D_{+}^{*}\)、\(T \in \mathrm{T}_{1}(n, \mathrm{K})\)，并且 D 和 T 连续依赖于 C。分解 X = UDT 的唯一性来自引理 10 的唯一性。

命题 7 中的同胚称为 \(\mathbf{GL}(n,\mathbf{K})\) 的 Iwasawa 分解。

群 \(\mathbf{G} = \mathbf{D}_{+}^{*} \cdot \mathbf{T}_{1}(n, \mathbf{K})\) 是 K 上对角元素为 \textgreater0 的上三角矩阵的集合。我们将该群的元素 \((z_{ij})\) 与元素

\[
\left(\left(z _ {i i}\right) _ {1 \leqslant i \leqslant n}, \left(z _ {i j}\right) _ {1 \leqslant i <   j \leqslant n}\right) \in \left(\mathbf {R} _ {+} ^ {*}\right) ^ {n} \times \mathrm{K} ^ {n (n - 1) / 2}.
\]

完全仿照例 4 的论证，当 \(\mathbf{K} = \mathbf{R}\) 时，可得该群上的右 Haar 测度为

\[
\left(\prod_ {i = 1} ^ {n} z _ {i i} ^ {- i}\right) \cdot \left(\bigotimes_ {i = 1} ^ {n} d z _ {i i}\right) \otimes \left(\bigotimes_ {i <   j} d z _ {i j}\right).
\]

然后应用 §2, No.~9 的 Prop. 13，可见当把 \(\mathbf{GL}(n,\mathbf{K})\) 与 \(\mathbf{U}(n,\mathbf{K})\times\mathbf{G}\) 通过映射 \((U,S)\mapsto US\) 认同后，\(\mathbf{GL}(n,\mathbf{K})\) 上的一个 Haar 测度为（当 \(\mathbf{K}=\mathbf{R}\) 时）

\[
\left(\prod_ {i = 1} ^ {n} z _ {i i} ^ {- i}\right) \cdot \alpha \otimes \left(\bigotimes_ {i = 1} ^ {n} d z _ {i i}\right) \otimes \left(\bigotimes_ {i <   n} d z _ {i j}\right),\tag{14}
\]

其中 \(\alpha\) 表示 \(\mathbf{U}(n,\mathbf{K})\) 上的一个 Haar 测度。

例 8. --- 厄米型空间。

在这个例子中，K 始终表示域 R、C、H 中的一个。记 \(\delta = \dim_{R} K\)（因而 \(\delta = 1, 2\) 或 4）。右向量空间上的一个厄米型 \(\Phi\)，其空间为 \(K^{n}\)，可写为

\[
\Phi (x, y) = \Phi (x _ {1}, \dots , x _ {n}, y _ {1}, \dots , y _ {n}) = \sum_ {i, j = 1} ^ {n} \overline {{{{x}}}} _ {i} h _ {i j} y _ {j}
\]

其中 \(h_{ij} = \overline{h}_{ji}\) 对所有 i 和 j 都成立。记 H 为 \(\mathbf{M}_{n}(\mathbf{K})\) 的厄米矩阵所成的 R 上向量空间。映射 \((h_{ij}) \mapsto \Phi\) 是 H 到 \(K^{n}\) 上的厄米型向量空间的同构，借此我们认同这两个空间。令 \(H_{+}^{*} \subset H\) 为 \(K^{n}\) 上非退化正厄米型的集合。集合 \(H_{+}^{*}\) 在 H 中是凸的；事实上，若 \(\Phi_{1}, \Phi_{2}\) 属于 \(H_{+}^{*}\)，且 \(\lambda, \mu\) 是满足 \textgreater0 和 \(\lambda + \mu = 1\) 的两个数，则 \(\lambda \Phi_{1} + \mu \Phi_{2}\) 显然是正厄米型；另一方面，若 \((\lambda \Phi_{1} + \mu \Phi_{2})(x, x) = 0\)，则 \(\Phi_{1}(x, x) = \Phi_{2}(x, x) = 0\)，因而 x = 0，所以 \(\lambda \Phi_{1} + \mu \Phi_{2}\) 是非退化的。现在来证明 \(H_{+}^{*}\) 是 H 的开子集。令 S 为满足 \(x = (x_{1}, \ldots, x_{n}) \in K^{n}\) 且 \(x_{1} \overline{x}_{1} + \cdots + x_{n} \overline{x}_{n} = 1\) 的 x 的集合；这是 \(K^{n}\) 的紧子集；若 \(\Phi \in H_{+}^{*}\)，函数 \(x \mapsto \Phi(x, x)\) 在 S 上连续且 \textgreater0，故其下确界 \textgreater0；若 \(\Phi' \in H\) 充分接近 \(\Phi\)，则 \(\Phi'(x, x) > 0\) 对所有 \(x \in S\) 成立，所以 \(\Phi'\) 是正的且非退化。

一般线性群 \(\mathbf{GL}(n,\mathbf{K})\) 在 \(\mathfrak{H}\) 上连续地右作用，即 \((X,\Phi)\mapsto\Phi\circ X\)，也就是按 \((X,H)\mapsto{}^{t}\overline{X}\cdot H\cdot X\) 作用，其中 \(H\) 表示与 \(\Phi\) 对应的厄米矩阵。显然，\(\mathfrak{H}_{+}^{*}\) 在 \(\mathbf{GL}(n,\mathbf{K})\) 下稳定。更确切地，根据 Alg., Ch. IX, §6, No.~1, Th. 1 的 Cor. 1，\(^{1}\) \(\mathfrak{H}_{+}^{*}\) 是 \(\mathbf{GL}(n,\mathbf{K})\) 作用在对应于型 \(\sum_{i=1}^{n}\overline{x}_{i}y_{i}\) 的单位矩阵 \(I_{n}\) 上所得的轨道。这个型的稳定子群是 \(\mathbf{U}(n,\mathbf{K})\)。根据附录 I 的引理 2，可以将 \(\mathfrak{H}_{+}^{*}\) 作为拓扑齐性空间与 \(\mathbf{GL}(n,\mathbf{K})/\mathbf{U}(n,\mathbf{K})\) 认同。

对于每个 \(X \in \mathbf{GL}(n, \mathbf{K})\)，令 \(\widetilde{X}\) 为自同构 \(H \mapsto {}^t \overline{X} \cdot H \cdot X\)，它作用在实向量空间 \(\mathfrak{H}\) 上。若 \(\mu\) 表示加法群 \(\mathfrak{H}\) 的 Haar 测度，则有 \(\widetilde{X}^{-1}(\mu) = |\det \widetilde{X}| \cdot \mu\)（§1, No.~10, Prop. 15 的 Cor. 1）。我们来证明

\[
| \det \widetilde {X} | = | N (X) | ^ {\lambda},\tag{15}
\]

其中 N 表示将 \(\mathbf{M}_{n}(\mathrm{K})\) 视为 R-代数时的范数，且 \(\lambda = 1 - \frac{\delta - 2}{\delta n}\)。只需验证 (15) 对于生成 \(\mathbf{GL}(n,\mathbf{K})\) 的一个生成系成立，因此（A, II, §10, No.~13, Prop. 14 的 Cor. 2）只需考虑以下三类 \(X\)：

\begin{enumerate}
\def\labelenumi{\alph{enumi})}
\tightlist
\item
  \(X\) 是如下形式的映射的矩阵
\end{enumerate}

\[
\left(x _ {1}, \dots , x _ {n}\right) \mapsto \left(x _ {\sigma (1)}, \dots , x _ {\sigma (n)}\right),
\]

其中 \(\sigma \in \mathfrak{S}_n\)。在此情形下，\(X\) 的某个幂等于 1，因此 \(|\det \widetilde{X}| = |\mathrm{N}(X)| = 1\)。

\begin{enumerate}
\def\labelenumi{\alph{enumi})}
\setcounter{enumi}{1}
\tightlist
\item
  \(X\) 是如下形式的映射的矩阵
\end{enumerate}

\[
\left(x _ {1}, \dots , x _ {n}\right) \mapsto \left(a x _ {1}, x _ {2}, \dots , x _ {n}\right).
\]

于是，若 \((h_{ij})\in \mathfrak{H}\)，则 \(\widetilde{X} ((h_{ij})) = (h_{ij}')\)，其中 \(h_{11}' = \overline{a} h_{11}a = |a|^2 h_{11}\)，\(h_{1i}' = \overline{a} h_{1i}\)（\(i > 1\)），并且 \(h_{ij}' = h_{ij}\)（\(i > 1\)，\(j > 1\)）；因此

\[
| \det \widetilde {X} | = | a | ^ {2} | a | ^ {\delta (n - 1)} = | a | ^ {2 + \delta (n - 1)}.
\]

另一方面，若 \(Y = (y_{ij}) \in \mathbf{M}_n(\mathbf{K})\)，则 \(XY = (y_{ij}')\)，其中 \(y_{1j}' = ay_{1j}\)，且 \(y_{ij}' = y_{ij}\)（\(i > 1\)）；因此 \(|\mathrm{N}(X)| = |a|^{\delta n}\)。公式 (15) 再次得到验证。

\begin{enumerate}
\def\labelenumi{\alph{enumi})}
\setcounter{enumi}{2}
\tightlist
\item
  \(X\) 是如下形式的映射的矩阵
\end{enumerate}

\[
\left(x _ {1}, \dots , x _ {n}\right) \mapsto \left(x _ {1} + b x _ {2}, x _ {2}, \dots , x _ {n}\right).
\]

于是 \(\widetilde{X}\big((h_{ij})\big) = (h_{ij}^{\prime})\)，其中 \(h_{11}^{\prime} = h_{11}\)，\(h_{12}^{\prime} = h_{12} + h_{11}b\)，\(h_{1i}^{\prime} = h_{1i}\)（\(i > 2\)），\(h_{22}^{\prime} = h_{22} + \overline{b}h_{12} + \overline{h}_{12}b + \overline{b}h_{11}b\)，\(h_{2i}^{\prime} = h_{2i} + \overline{b}h_{1i}\)（\(i > 2\)），\(h_{ij}^{\prime} = h_{ij}\)（\(i > 2\)，\(j > 2\)）。考虑引理 6，可见 \(|\det \widetilde{X}| = 1\)。同样可验证 \(|\mathrm{N}(X)| = 1\)，这就完成了公式 (15) 的证明。

由此，H 上的测度 \(|\mathrm{N}(H)|^{-\lambda/2}d\mu(H)\) 在 \(\mathbf{GL}(n,\mathbf{K})\) 下不变，因为

\[
\begin{array}{r l} & {\widetilde {X} ^ {- 1} \big (| \mathrm{N} (H) | ^ {- \lambda / 2} d \mu (H) \big) = \big | \mathrm{N} \big (\widetilde {X} (H) \big) \big | ^ {- \lambda / 2} \cdot | \det \widetilde {X} | d \mu (H)} \\ & {\qquad = | \mathrm{N} (H) | ^ {- \lambda / 2} | \mathrm{N} (X) | ^ {- \lambda} | \det \widetilde {X} | d \mu (H) = | \mathrm{N} (H) | ^ {- \lambda / 2} d \mu (H).} \end{array}
\]

若 \(H \in \mathfrak{H}_+^*\)，则 \(H = \widetilde{X}(\mathrm{I}_n) = {}^t\overline{X}X\)，其中 \(X \in \mathbf{GL}(n,\mathbf{K})\)，因此 \(\mathrm{N}(H) = \overline{\mathrm{N}(X)}\mathrm{N}(X) > 0\)。所以在 \(\mathfrak{H}_+^*\) 上，唯一（差一个常数因子）的、在 \(\mathbf{GL}(n,\mathbf{K})\) 下不变的测度（cf.~§2, No.~6, Th. 3）是测度

\[
d \gamma (H) = \mathrm{N} (H) ^ {- \lambda / 2} d \mu (H).
\]

特别地，

\[
d \gamma (H) = (\det H) ^ {- (n + 1) / 2} d \mu (H) \quad \text { when } K = R,
\]

\[
d \gamma (H) = (\det H) ^ {- n} d \mu (H) \quad \text { when } \mathrm{K} = \mathrm{C}.
\]

